%! Graphing Rational Functions — Unit 4, Lesson 4.5 — Group Activity (Answer Key)
\documentclass[10pt]{article}
\usepackage{algebra2-article}
\usepackage{algebra2-key}   % pulls in -boxes; do NOT also load -boxes
\newcommand{\TallMath}[1]{$\displaystyle #1\rule[-1.4em]{0pt}{3.2em}$}
% Standard coordinate grid + axes: {xmin}{xmax}{ymin}{ymax}
\newcommand{\lgrid}[4]{%
  \draw[step=1, linegray!40, very thin] (#1,#3) grid (#2,#4);
  \draw[->, semithick, charcoal] (#1,0)--(#2,0) node[below right, font=\scriptsize]{$x$};
  \draw[->, semithick, charcoal] (0,#3)--(0,#4) node[left, font=\scriptsize]{$y$};}

\begin{document}

\pageheader{Unit 4, Lesson 4.5}{Group Activity --- Answer Key}
\namepartnerperiod

\vspace{0.08in}

\begin{headlinebox}{forestbg}
\small\textbf{Assemble the Graph.} Same five steps every time, and always in this order:
\textbf{(1)} factor everything; \textbf{(2)} restrictions off the \emph{original} denominator, then the
cancel test --- hole or wall; \textbf{(3)} degree comparison for the horizontal asymptote;
\textbf{(4)} intercepts, \emph{computed}; \textbf{(5)} sign chart, with the $x$-intercepts and the
vertical asymptotes as the boundary points. Nothing gets read off the face of the equation today.
\end{headlinebox}

\vspace{0.06in}

\begin{tcolorbox}[colback=white, colframe=black!40, title=\textbf{Tier R --- Remediate},
    fonttitle=\bfseries, arc=1.5mm, left=3mm, right=3mm, top=2mm, bottom=2mm, breakable]
\textbf{Part 1 --- Hole or wall?} Factor if you need to, list the restrictions off the \emph{original}
denominator, then sort each one.
\begin{center}
\renewcommand{\arraystretch}{1.7}
\begin{tabularx}{\linewidth}{@{}l X X X@{}}
\hline
\textbf{Function} & \textbf{Restrictions} & \textbf{Hole at\dots} & \textbf{Vertical asymptote(s)} \\
\hline
\TallMath{\frac{x-1}{(x-1)(x+4)}}  & \ans{$x\neq1,-4$} & \ans{$x=1$} & \ans{$x=-4$} \\
\TallMath{\frac{x+2}{(x-3)(x+5)}}  & \ans{$x\neq3,-5$} & \ans{none} & \ans{$x=3$, $x=-5$} \\
\TallMath{\frac{x^2-9}{(x+3)(x-1)}} & \ans{$x\neq-3,1$} & \ans{$x=-3$} & \ans{$x=1$} \\
\TallMath{\frac{5}{(x-2)(x+2)}}    & \ans{$x\neq2,-2$} & \ans{none} & \ans{$x=2$, $x=-2$} \\
\hline
\end{tabularx}
\end{center}

\vspace{4pt}
\textbf{Part 2 --- Horizontal asymptote, by degrees only.} Write $y=$ something, or ``none.''
\begin{enumerate}[label=(\alph*), leftmargin=2.1em, itemsep=4pt]
  \item $y=\dfrac{x+1}{x^2-9}$ \hfill top deg \ans{$1$} \; bottom deg \ans{$2$} \;
        HA \ans{$y=0$}
  \item $y=\dfrac{4x^2+1}{2x^2-x}$ \hfill top deg \ans{$2$} \; bottom deg \ans{$2$} \;
        HA \ans{$y=2$}
  \item $y=\dfrac{x^2-4}{x-1}$ \hfill top deg \ans{$2$} \; bottom deg \ans{$1$} \;
        HA \ans{none}
  \item $y=\dfrac{3x}{x+7}$ \hfill top deg \ans{$1$} \; bottom deg \ans{$1$} \;
        HA \ans{$y=3$}
\end{enumerate}

\vspace{4pt}
\textbf{Part 3 --- One full build, step by step.} \; $f(x)=\dfrac{3}{(x-1)(x+3)}$
\begin{enumerate}[label=\arabic*., leftmargin=2.1em, itemsep=5pt]
  \item Restrictions: $x\neq$ \ans{$1$}, $x\neq$ \ans{$-3$}. Anything cancel? \ans{no}
        \; So the vertical asymptotes are \ans{$x=1$} and \ans{$x=-3$}.
  \item Degrees: top \ans{$0$} (the numerator is a constant), bottom \ans{$2$}, so the
        horizontal asymptote is \ans{$y=0$}.
  \item $x$-intercept: set the numerator to zero. What happens? \ans{$3=0$ is impossible --- there is no $x$-intercept}
  \item $y$-intercept: $f(0)=\dfrac{3}{(-1)(3)}=$ \ans{$-1$}, so the point is \ans{$(0,-1)$}.
  \item Sign chart. With no $x$-intercept, the only boundary points are the two \ans{asymptotes}.
        \begin{center}
        \renewcommand{\arraystretch}{1.5}
        \begin{tabularx}{0.9\linewidth}{|C{2.6cm}|C{1.6cm}|C{2.0cm}|X|}
        \hline
        \rowcolor{forestbg}
        \textbf{Interval} & \textbf{Test $x$} & \textbf{Sign} & \textbf{Above / below} \\ \hline
        $x<-3$   & $-4$ & \ans{$+$} & \ans{above} \\ \hline
        $-3<x<1$ & $0$  & \ans{$-$} & \ans{below} \\ \hline
        $x>1$    & $2$  & \ans{$+$} & \ans{above} \\ \hline
        \end{tabularx}
        \end{center}
  \item Domain: \ans{$(-\infty,-3)\cup(-3,1)\cup(1,\infty)$}
  \item This graph never crosses its horizontal asymptote. Why not? \ansline{crossing $y=0$ would need $f(x)=0$, and a fraction is zero only when its numerator is --- this numerator is the constant $3$}
\end{enumerate}
\end{tcolorbox}

\begin{tcolorbox}[colback=white, colframe=black!40, title=\textbf{Tier A --- Approaching Proficiency},
    fonttitle=\bfseries, arc=1.5mm, left=3mm, right=3mm, top=2mm, bottom=2mm, breakable]
\textbf{Part 1 --- Build it, then check yourself against the real graph.} \;
$f(x)=\dfrac{x-2}{x^2-x-12}$
\begin{enumerate}[label=\arabic*., leftmargin=2.1em, itemsep=4pt]
  \item Factored: $\dfrac{x-2}{(\text{\ans{$x-4$}})(\text{\ans{$x+3$}})}$ \quad Restrictions:
        $x\neq$ \ans{$4$}, $x\neq$ \ans{$-3$} \quad Holes: \ans{none}
  \item Vertical asymptotes \ans{$x=4$} and \ans{$x=-3$}. \quad Horizontal asymptote
        \ans{$y=0$}.
  \item $x$-intercept \ans{$(2,0)$}. \quad $y$-intercept \ans{$(0,\frac16)$}.
  \item Sign chart --- boundary points \ans{$-3$, $2$, $4$}.
        \begin{center}
        \renewcommand{\arraystretch}{1.5}
        \begin{tabularx}{0.92\linewidth}{|C{2.4cm}|C{1.5cm}|C{1.8cm}|X|}
        \hline
        \rowcolor{forestbg}
        \textbf{Interval} & \textbf{Test $x$} & \textbf{Sign} & \textbf{Above / below} \\ \hline
        $x<-3$   & $-4$ & \ans{$-$} & \ans{below} \\ \hline
        $-3<x<2$ & $0$  & \ans{$+$} & \ans{above} \\ \hline
        $2<x<4$  & $3$  & \ans{$-$} & \ans{below} \\ \hline
        $x>4$    & $5$  & \ans{$+$} & \ans{above} \\ \hline
        \end{tabularx}
        \end{center}
  \item Now compare with the graph below. Does your chart match it? \ans{yes} \\[3pt]
        The graph \emph{crosses} its horizontal asymptote $y=0$ at \ans{$(2,0)$}. Explain in one
        sentence why a crossing was guaranteed here but impossible for every function in Lesson 4.4.
        \ansline{this numerator has a zero, so $f(x)=0$ has a solution; in 4.4 the numerator was a nonzero constant, so it never could}
\end{enumerate}
\begin{center}
\begin{tikzpicture}[scale=0.36]
  \lgrid{-8}{9}{-4}{4}
  \draw[navy, dashed, semithick] (-3,-4)--(-3,4);
  \draw[navy, dashed, semithick] (4,-4)--(4,4);
  \node[navy, font=\scriptsize, above left] at (-3,3.2) {$x=-3$};
  \node[navy, font=\scriptsize, above right] at (4,3.2) {$x=4$};
  \begin{scope}
    \clip (-8,-4) rectangle (9,4);
    \draw[very thick, forest, domain=-8:-3.2, samples=120, smooth] plot (\x,{((\x)-2)/(((\x)-4)*((\x)+3))});
    \draw[very thick, forest, domain=-2.8:3.93, samples=150, smooth] plot (\x,{((\x)-2)/(((\x)-4)*((\x)+3))});
    \draw[very thick, forest, domain=4.07:9, samples=120, smooth] plot (\x,{((\x)-2)/(((\x)-4)*((\x)+3))});
  \end{scope}
  \fill[charcoal] (2,0) circle (3.2pt) node[below left, font=\scriptsize]{$(2,0)$};
\end{tikzpicture}
\end{center}

\vspace{4pt}
\textbf{Part 2 --- Build one with a hole.} \; $g(x)=\dfrac{x^2-2x-3}{x^2-9}$
\begin{enumerate}[label=\arabic*., leftmargin=2.1em, itemsep=4pt]
  \item Factored: $\dfrac{(\text{\ans{$x-3$}})(\text{\ans{$x+1$}})}%
        {(\text{\ans{$x-3$}})(\text{\ans{$x+3$}})}$ \quad Restrictions: $x\neq$ \ans{$3$},
        $x\neq$ \ans{$-3$}
  \item The factor \ans{$x-3$} cancels, so there is a hole at $x=$ \ans{$3$}. Simplified:
        \ans{$\frac{x+1}{x+3}$}. Hole coordinates: \ans{$(3,\frac23)$}
  \item Vertical asymptote \ans{$x=-3$}. \quad Horizontal asymptote \ans{$y=1$} (degrees
        \ans{$2$} and \ans{$2$}).
  \item $x$-intercept \ans{$(-1,0)$}. \quad $y$-intercept \ans{$(0,\frac13)$}. \\[3pt]
        The \emph{original} numerator is also zero at $x=3$. Why is $(3,0)$ \textbf{not} an
        $x$-intercept? \ans{$x=3$ is not in the domain --- it is the hole, so the graph has no point there at all}
  \item Sign chart --- boundary points \ans{$-3$ and $-1$}. \; Test $x=-4$: \ans{$+$} \;
        $x=-2$: \ans{$-$} \; $x=0$: \ans{$+$}
  \item Domain \ans{$(-\infty,-3)\cup(-3,3)\cup(3,\infty)$}
\end{enumerate}

\vspace{4pt}
\textbf{Part 3 --- Error analysis.} A student analyzed $h(x)=\dfrac{x^2-16}{x^2-3x-4}$ and wrote:
\begin{center}
\fbox{\parbox{0.80\linewidth}{\centering \vspace{2pt}
``Vertical asymptotes at $x=4$ and $x=-1$. \\
$x$-intercepts at $x=4$ and $x=-4$. \\
Horizontal asymptote $y=1$.''\vspace{2pt}}}
\end{center}
\begin{enumerate}[label=(\alph*), leftmargin=2.1em, itemsep=5pt]
  \item Factor it: $h(x)=\dfrac{(\text{\ans{$x-4$}})(\text{\ans{$x+4$}})}%
        {(\text{\ans{$x-4$}})(\text{\ans{$x+1$}})}$
  \item \textbf{One} of the student's three claims is correct. Which? \ans{the HA, $y=1$}
  \item Fix the vertical asymptotes: \ans{only $x=-1$}, because $x=4$ is really a \ans{hole},
        located at \ans{$(4,\frac85)$}.
  \item Fix the $x$-intercepts: \ans{only $(-4,0)$}. In one sentence, state the rule the student broke.
        \ansline{the cancel test comes first --- a factor that cancels is a hole, not a wall, and its zero is not an intercept}
\end{enumerate}
\end{tcolorbox}

\begin{tcolorbox}[colback=white, colframe=black!40, title=\textbf{Tier E --- Extension},
    fonttitle=\bfseries, arc=1.5mm, left=3mm, right=3mm, top=2mm, bottom=2mm, breakable]
\textbf{Part 1 --- Build one backwards.} Write the equation of a rational function that has a hole at
$x=1$, a vertical asymptote at $x=-4$, a horizontal asymptote $y=2$, and an $x$-intercept at $(3,0)$.
\begin{center}
$y=$ \ans{$\dfrac{2(x-3)(x-1)}{(x+4)(x-1)}$}
\end{center}
\begin{enumerate}[label=(\alph*), leftmargin=2.1em, itemsep=5pt]
  \item Explain how each of the four requirements forced part of your answer.
        \ansline{hole at $1$ $\Rightarrow$ the factor $x-1$ must appear on \emph{both} top and bottom; VA at $-4$ $\Rightarrow$ $x+4$ on the bottom only; $x$-intercept $3$ $\Rightarrow$ $x-3$ on the top only; HA $y=2$ $\Rightarrow$ the degrees already match ($2$ and $2$), so the leading coefficients must be in the ratio $2:1$}
  \item Where exactly is the hole? \ans{$(1,-\frac45)$}
  \item Could the $x$-intercept have been at $(1,0)$ instead? Explain. \ansline{no --- $x=1$ is the hole, so it is not in the domain, and an intercept has to be a point the graph actually reaches}
\end{enumerate}

\vspace{5pt}
\textbf{Part 2 --- Enrichment: what happens when the top outranks the bottom?}
\emph{(This one is for curiosity. It will not be on the test.)}

\vspace{2pt}
$f(x)=\dfrac{x^2-4}{x+1}$ has top degree $2$ and bottom degree $1$, so by the degree comparison it has
\ans{no} horizontal asymptote. Divide anyway --- Lesson 3.3, long division:
\begin{center}
$\dfrac{x^2-4}{x+1}=$ \ans{$x-1$} $+\dfrac{\text{\ans{$-3$}}}{x+1}$
\end{center}
\begin{center}
\begin{tikzpicture}[scale=0.30]
  \lgrid{-8}{7}{-8}{8}
  \draw[navy, dashed, semithick] (-1,-8)--(-1,8);
  \node[navy, font=\scriptsize, above right] at (-1,7.2) {$x=-1$};
  \begin{scope}
    \clip (-8,-8) rectangle (7,8);
    \draw[redacc, dashed, semithick] (-8,-9)--(7,6);
    \draw[very thick, forest, domain=-8:-1.3, samples=140, smooth] plot (\x,{((\x)*(\x)-4)/((\x)+1)});
    \draw[very thick, forest, domain=-0.62:7, samples=140, smooth] plot (\x,{((\x)*(\x)-4)/((\x)+1)});
  \end{scope}
  \node[redacc, font=\scriptsize, below right] at (5.2,4.2) {$y=x-1$};
  \fill[charcoal] (2,0) circle (3.4pt);
  \fill[charcoal] (-2,0) circle (3.4pt);
\end{tikzpicture}
\end{center}
\vspace{-4pt}
\begin{enumerate}[label=(\alph*), leftmargin=2.1em, itemsep=5pt]
  \item Far from the origin, the fraction on the right dies to \ans{$0$}, so the graph runs
        alongside the \emph{line} $y=$ \ans{$x-1$} --- the dashed line in the figure. That line is
        called a \textbf{slant} (or oblique) \textbf{asymptote}.
  \item Vertical asymptote \ans{$x=-1$}. \; $x$-intercepts \ans{$(2,0)$ and $(-2,0)$}. \; $y$-intercept
        \ans{$(0,-4)$}.
  \item Compare with Lesson 4.4's disguise, $\frac{x-1}{x-3}=1+\frac{2}{x-3}$. What is the same about
        the two rewrites, and what is different about the leftover? \ansline{both divide and keep the remainder as a fraction; in 4.4 the quotient was a \emph{constant}, so the graph settled onto a horizontal line --- here the quotient is \emph{linear}, so it settles onto a slanted one}
\end{enumerate}

\vspace{4pt}
\textbf{Part 3 --- The round trip.} You ride $30$ miles out at a steady $x$ miles per hour, then ride
the same $30$ miles back at $x+10$ miles per hour (downhill, or with the wind). Total time, in hours:
\[ T(x)=\frac{30}{x}+\frac{30}{x+10}. \]
\begin{enumerate}[label=(\alph*), leftmargin=2.1em, itemsep=5pt]
  \item Combine into a single fraction (Lesson 4.3): \; $T(x)=$ \ans{$\dfrac{60x+300}{x(x+10)}$}
  \item Now run the build on it. Restrictions \ans{$x\neq0,-10$}; vertical asymptotes \ans{$x=0$, $x=-10$};
        horizontal asymptote \ans{$y=0$}; $x$-intercept \ans{$(-5,0)$}.
  \item Two of those four features are meaningless for a bicycle. Circle the two that \emph{do} mean
        something, and say what each one means.
        \ansline{$x=0$ and $y=0$. The wall at $x=0$ says that as your speed drops toward zero the trip time grows without bound. The floor at $y=0$ says the faster you ride the closer the total time gets to zero --- approached, never reached. The wall at $x=-10$ and the intercept $(-5,0)$ are negative speeds: real algebra, no bicycle.}
  \item $T(10)=$ \ans{$4.5$} h, \; $T(20)=$ \ans{$2.5$} h, \; $T(30)=$ \ans{$1.75$} h, \;
        $T(40)=$ \ans{$1.35$} h.
  \item Going from $10$ to $20$ mph saved \ans{$2$} hours. Going from $20$ to $40$ mph saved
        \ans{$1.15$} hours --- less, even though you doubled your speed both times. Explain what
        the graph is doing near its horizontal asymptote, and why no amount of speed gets the trip
        time to zero. \ansline{the graph is flattening onto $y=0$, so each extra mile per hour buys less time than the one before; $T(x)=0$ would need $60x+300=0$, i.e.\ a negative speed, so the time is always positive}
\end{enumerate}
\end{tcolorbox}

\begin{teachernote}
\textbf{Tier R} Part 3: the numerator is the constant $3$, so there is \emph{no} $x$-intercept.
\textbf{Tier A} Part 2 item 4 is the point of the tier --- the original numerator vanishes at $x=3$ and
$(3,0)$ is still not an intercept, because $3$ is not in the domain. In Part 3 the only surviving claim
is the HA, the one feature the cancel test does not touch. \textbf{Tier E} Part 1 is the best
diagnostic here; accept any equivalent form. \textbf{Part 2 is enrichment only: slant asymptotes are
not in A2.F.2h and are never assessed} --- say so out loud. Part 3 reverses 4.4's dilution model: there
the \emph{vertical} asymptote was meaningless, here $x=0$ carries the meaning.
\end{teachernote}

\end{document}
